\section{Exactness and points}
\label{section-exactness-points}

\noindent
In any category we have the notion of epimorphism, monomorphism,
isomorphism, etc.

\begin{lemma}
\label{lemma-points-exactness}
Let $X$ be a topological space. Let $\varphi : \mathcal{F} \to \mathcal{G}$
be a morphism of sheaves of sets on $X$.
\begin{enumerate}
\item The map $\varphi$ is a monomorphism in the category of sheaves
if and only if for all $x \in X$ the map
$\varphi_x : \mathcal{F}_x \to \mathcal{G}_x$
is injective.
\item The map $\varphi$ is an epimorphism in the category of sheaves
if and only if for all $x \in X$ the map
$\varphi_x : \mathcal{F}_x \to \mathcal{G}_x$
is surjective.
\item The map $\varphi$ is an isomorphism in the category of sheaves
if and only if for all $x \in X$ the map
$\varphi_x : \mathcal{F}_x \to \mathcal{G}_x$
is bijective.
\end{enumerate}
\end{lemma}

\begin{proof}
Omitted.
\end{proof}

\noindent
It follows that in the category of sheaves of sets
the notions epimorphism and monomorphism can be described
as follows.

\begin{definition}
\label{definition-injective-surjective}
Let $X$ be a topological space.
\begin{enumerate}
\item A presheaf $\mathcal{F}$ is called a {\it subpresheaf} of a presheaf
$\mathcal{G}$ if $\mathcal{F}(U) \subset \mathcal{G}(U)$ for all open
$U \subset X$ such that the restriction maps of $\mathcal{G}$ induce the
restriction maps of $\mathcal{F}$. If $\mathcal{F}$ and
$\mathcal{G}$ are sheaves, then $\mathcal{F}$ is called a {\it subsheaf}
of $\mathcal{G}$. We sometimes indicate this by the notation
$\mathcal{F} \subset \mathcal{G}$.
\item A morphism of presheaves of sets $\varphi : \mathcal{F} \to \mathcal{G}$
on $X$ is called {\it injective} if and only if
$\mathcal{F}(U) \to \mathcal{G}(U)$ is injective for all $U$ open in $X$.
\item A morphism of presheaves of sets $\varphi : \mathcal{F} \to \mathcal{G}$
on $X$ is called {\it surjective} if and only if
$\mathcal{F}(U) \to \mathcal{G}(U)$ is surjective for all $U$ open in $X$.
\item A morphism of sheaves of sets $\varphi : \mathcal{F} \to \mathcal{G}$
on $X$ is called {\it injective} if and only if
$\mathcal{F}(U) \to \mathcal{G}(U)$ is injective for all $U$ open in $X$.
\item A morphism of sheaves of sets $\varphi : \mathcal{F} \to \mathcal{G}$
on $X$ is called {\it surjective} if and only if for every open
$U$ of $X$ and every section $s$ of $\mathcal{G}(U)$ there exists an
open covering $U = \bigcup U_i$ such that $s|_{U_i}$ is in
the image of $\mathcal{F}(U_i) \to \mathcal{G}(U_i)$ for all $i$.
\end{enumerate}
\end{definition}


\begin{lemma}
\label{lemma-characterize-epi-mono}
Let $X$ be a topological space.
\begin{enumerate}
\item Epimorphisms (resp.\ monomorphisms) in the category of
presheaves are exactly the surjective (resp.\ injective) maps
of presheaves.
\item Epimorphisms (resp.\ monomorphisms) in the category of
sheaves are exactly the surjective (resp.\ injective) maps
of sheaves, and are exactly those maps which are surjective
(resp.\ injective) on all the stalks.
\end{enumerate}
\end{lemma}

\begin{proof}
Omitted.
\end{proof}


\begin{lemma}
\label{lemma-check-homomorphism-stalks}
let $X$ be a topological space.
Let $(\mathcal{C}, F)$ be a type of algebraic structure.
Suppose that $\mathcal{F}$, $\mathcal{G}$ are sheaves on $X$
with values in $\mathcal{C}$.
Let $\varphi : \mathcal{F} \to \mathcal{G}$
be a map of the underlying sheaves of sets.
If for all points $x \in X$ the map
$\mathcal{F}_x \to \mathcal{G}_x$
is a morphism of algebraic structures,
then $\varphi$ is a morphism of sheaves of algebraic structures.
\end{lemma}

\begin{proof}
Let $U$ be an open subset of $X$. Consider the diagram of
(underlying) sets
$$
\xymatrix{
\mathcal{F}(U) \ar[r] \ar[d] &
\prod_{x \in U} \mathcal{F}_x \ar[d] \\
\mathcal{G}(U) \ar[r] &
\prod_{x \in U} \mathcal{G}_x
}
$$
By assumption, and previous results, all but the left vertical
arrow are morphisms of algebraic structures. In addition the
bottom horizontal arrow is injective,
see Lemma \ref{lemma-sheaf-subset-stalks}.
Hence we conclude by Lemma \ref{lemma-image-contained-in},
see also Example \ref{example-application-lemma-image-contained-in}
\end{proof}

\noindent
Short exact sequences of abelian sheaves, etc will be discussed
in the chapter on sheaves of modules. See
Modules, Section \ref{modules-section-kernels}.
